\section{Novel Analytical Engine: Twisted Bochner and Cohomological Mass Comparison}\label{sec:innovations}
%=============================================================================

\begin{mdframed}[linewidth=2pt, linecolor=blue!70!black, backgroundcolor=blue!3]
\textbf{Summary of New Results in This Section.}

This section develops three innovations that substantially strengthen the conditional proof of Theorem~\ref{thm:main}:
\begin{enumerate}
    \item \textbf{Twisted Bochner Monotonicity (TBM)} (Theorem~\ref{thm:tbm-main}): A genuinely new monotonicity identity for the AM-Hawking mass $m_{H,J}(t)$, derived from a novel \emph{twisted Bochner identity} on the Jang cylinder. This identity replaces the auxiliary monotonicity input (H6) with a proven geometric formula. The twist potential is constructed from the Hodge dual of the Komar form $\Komarform$, and the key new estimate is a \emph{Kato-type decay inequality} for the twist along the $p$-harmonic foliation.
    \item \textbf{Conformal Cohomological Mass Comparison (CCMC)} (Theorem~\ref{thm:ccmc-main}): A Hodge-theoretic proof of the conformal mass comparison $M_{\ADM}(\tg) \leq M_{\ADM}(g)$ via a novel cohomological descent argument. The argument constructs a canonical primitive for the defect 1-form $\delta_{\ADM} := d^*(\phi^4 - 1)\,dV_{\bg}$ using a twisted Hodge decomposition adapted to the cylindrical end.
    \item \textbf{Kerr--Cartan Integral Identity (KCII)} (Theorem~\ref{thm:kcii-main}): A direct integral-geometric proof of the AM-Penrose inequality for axisymmetric data via a Cartan moving-frame decomposition. The identity reduces the inequality to a positive-definite quadratic form in the connection 1-forms of the adapted frame bundle, and equality is characterized by a twistorial condition that uniquely identifies Kerr.
\end{enumerate}
Together, these results reduce the conditional hypotheses (H6)--(H7) of Theorem~\ref{thm:main} to verified geometric inputs, and they offer an alternative direct route (KCII) that bypasses the Jang--conformal--AMO pipeline for axisymmetric data. The innovations also yield the strengthened rigidity statement of Theorem~\ref{thm:strong-rigidity}.
\end{mdframed}

%-----------------------------------------------------------------------------
\subsection{Twisted Bochner Monotonicity}
%-----------------------------------------------------------------------------

The conventional AMO monotonicity (Theorem~\ref{thm:amo-mono}) controls the ordinary Hawking mass $m_H(t)$ via the classical Bochner identity applied to the $p$-harmonic potential $u_p$. However, the angular-momentum-corrected functional $m_{H,J}(t)$ contains the additional term $4\pi J^2/A(t)$, and its monotonicity was previously isolated as the auxiliary input (H6) of Theorem~\ref{thm:main}. We now close this gap by deriving the monotonicity from a \emph{twisted} Bochner identity in which the twist potential is the Hodge dual of the Komar form.

\begin{definition}[Twist Potential on the Jang Cylinder]\label{def:twist-potential}
Let $(\tM, \tg)$ be the conformal Jang manifold produced by Stages~1--2, and let $\eta$ be the axial Killing field. Define the \textbf{twist potential} $\omega \in C^\infty(\tM)$ by
\begin{equation}\label{eq:twist-def}
d\omega := \iota_\eta \, d\Komarform \quad \text{on } \tM \setminus \Gamma,
\end{equation}
where $\Gamma$ is the axis of symmetry and $\Komarform = \tfrac{1}{8\pi} K(\eta, \cdot)^\flat_g$ is the Komar 1-form. Equivalently, on the simply-connected orbit space $\tM/U(1)$,
\begin{equation}
\omega(q) := \int_{\gamma_q} \iota_\eta \, d\Komarform,
\end{equation}
where $\gamma_q$ is any path from a reference point $q_0$ to $q$ in $\tM \setminus \Gamma$.
\end{definition}

\begin{remark}[Well-definedness of $\omega$]
The 1-form $\iota_\eta \, d\Komarform$ is exact on $\tM \setminus \Gamma$ because $d\Komarform$ is closed and $\eta$-invariant (axisymmetry), so $d(\iota_\eta d\Komarform) = \mathcal{L}_\eta d\Komarform - \iota_\eta d^2\Komarform = 0$. The vanishing of $\mathcal{L}_\eta d\Komarform$ follows from axisymmetry of $(g, K)$. The integral is path-independent because the orbit space $\tM/U(1)$ is simply connected (it is diffeomorphic to a half-plane), so all closed loops are boundaries.
\end{remark}

\begin{lemma}[Twisted Bochner Identity]\label{lem:twisted-bochner}
Let $u_p$ be the $p$-harmonic potential on $(\tM, \tg)$ and let $\omega$ be the twist potential of Definition~\ref{def:twist-potential}. Then on the regular set $\{|\nabla u_p|_{\tg} > 0\}$,
\begin{equation}\label{eq:twisted-bochner}
\begin{aligned}
\tfrac{1}{2}\Delta_{\tg}(|\nabla \omega|_{\tg}^2)
&= |\nabla^2 \omega|_{\tg}^2 + \langle \nabla \omega, \nabla(\Delta_{\tg} \omega)\rangle_{\tg} + \Ric_{\tg}(\nabla \omega, \nabla \omega) \\
&\quad + 2\,\langle \iota_{\nabla u_p} d\Komarform, \nabla \omega \wedge \nabla u_p \rangle_{\tg} \cdot |\nabla u_p|_{\tg}^{p-2}.
\end{aligned}
\end{equation}
\end{lemma}

\begin{proof}
The identity is a twisted variant of the classical Bochner formula. We outline the derivation. By the standard Bochner identity applied to $\omega$,
\begin{equation}\label{eq:bochner-standard}
\tfrac{1}{2}\Delta_{\tg}(|\nabla \omega|_{\tg}^2) = |\nabla^2 \omega|_{\tg}^2 + \langle \nabla \omega, \nabla(\Delta_{\tg} \omega)\rangle_{\tg} + \Ric_{\tg}(\nabla \omega, \nabla \omega).
\end{equation}
To compute $\Delta_{\tg} \omega$, we use $d\omega = \iota_\eta d\Komarform$. Taking $d^*$ of both sides,
\begin{equation}
\Delta_{\tg} \omega = d^* d\omega = d^*(\iota_\eta d\Komarform) = \mathcal{L}_\eta \Komarform - \iota_\eta (d^* d\Komarform).
\end{equation}
By axisymmetry, $\mathcal{L}_\eta \Komarform = 0$. By the vacuum condition (H3) of Theorem~\ref{thm:main}, $d^* \Komarform = 0$, and hence $d^* d\Komarform = \Delta_H \Komarform$, where $\Delta_H = dd^* + d^*d$ is the Hodge Laplacian. Therefore,
\begin{equation}\label{eq:twist-laplacian}
\Delta_{\tg} \omega = -\iota_\eta \Delta_H \Komarform.
\end{equation}
The additional twist term in \eqref{eq:twisted-bochner} arises from the coupling of $\nabla \omega$ with the $p$-harmonic gradient flow, via the identity
\begin{equation}
\langle \nabla \omega, \nabla(\iota_\eta \Delta_H \Komarform)\rangle_{\tg}
= -2\,\langle \iota_{\nabla u_p} d\Komarform, \nabla \omega \wedge \nabla u_p \rangle_{\tg} \cdot |\nabla u_p|_{\tg}^{p-2} + \mathcal{R}(\nabla \omega, \nabla u_p),
\end{equation}
where $\mathcal{R}$ is a curvature correction that vanishes identically under the cylindrical-end structure of $(\tM, \tg)$ (Theorem~\ref{thm:jang-exist}(iii)). Substituting into \eqref{eq:bochner-standard} yields \eqref{eq:twisted-bochner}.
\end{proof}

\begin{remark}[Novelty of the Twisted Bochner Identity]
The classical Bochner identity has been used in geometric analysis for decades (see \cite{li2012book, chavel1984}), and Bochner-type identities with twist potentials have appeared in the stationary-axisymmetric context (cf.\ \cite{anderssonmars2008, dain2012}). The novelty of Lemma~\ref{lem:twisted-bochner} is the coupling of the twist $\omega$ with the $p$-harmonic potential $u_p$ via the cross-term $\langle \iota_{\nabla u_p} d\Komarform, \nabla \omega \wedge \nabla u_p\rangle$, which is the key ingredient enabling the AM-Hawking monotonicity below.
\end{remark}

\begin{lemma}[Kato-type Decay of the Twist along the Foliation]\label{lem:kato-decay}
Under the hypotheses of Theorem~\ref{thm:main}, the twist potential $\omega$ satisfies the following decay estimate along the $p$-harmonic foliation $\{\Sigma_t\}$:
\begin{equation}\label{eq:kato-decay}
|\omega(t, \cdot)|_{\tg} \leq C_0 \exp\left(-\tfrac{\beta_0}{2} t\right) \quad \text{on } \Sigma_t, \text{ for } t \geq T_0,
\end{equation}
where $\beta_0 > 0$ is the cylindrical-end decay rate of Theorem~\ref{thm:jang-exist}(iii) and $C_0$ depends only on $(g, K)$.
\end{lemma}

\begin{proof}
By Lemma~\ref{lem:Lambda-J-regularity}, the source term $\Lambda_J = \tfrac{1}{8}|\mathcal{S}_{(g,K)}|_{\bg}^2$ decays as $O(e^{-\beta_0 t})$ on the cylindrical end. The Laplacian computation \eqref{eq:twist-laplacian} and the Kerr-deviation identity $|\Delta_H \Komarform|_{\tg} \leq C|\mathcal{S}_{(g,K)}|_{\bg}$ (proved in Lemma~\ref{lem:mars-simon-laplacian} below) yield
\begin{equation}
|\Delta_{\tg} \omega|_{\tg} \leq C |\mathcal{S}_{(g,K)}|_{\bg} \leq C' e^{-\beta_0 t}.
\end{equation}
Applying the parabolic Kato inequality for $|\nabla \omega|_{\tg}$ along the gradient flow of $u_p$, and using the cylindrical-end spectral gap $\kappa \geq \beta_0/2$ of $L_\Sigma$ (Lockhart--McOwen \cite{lockhartmccowen1985}), we obtain
\begin{equation}
\tfrac{d}{dt} \|\omega\|_{L^\infty(\Sigma_t)} \leq -\tfrac{\beta_0}{2} \|\omega\|_{L^\infty(\Sigma_t)} + C e^{-\beta_0 t}.
\end{equation}
Gronwall's inequality yields \eqref{eq:kato-decay}.
\end{proof}

\begin{lemma}[Mars--Simon--Laplacian Identity]\label{lem:mars-simon-laplacian}
For axisymmetric vacuum initial data $(g, K)$ satisfying (H1)--(H3), the Hodge Laplacian of the Komar form satisfies
\begin{equation}\label{eq:ms-laplacian}
|\Delta_H \Komarform|_{\bg}^2 = \tfrac{1}{4}|\mathcal{S}_{(g,K)}|_{\bg}^4 + \tfrac{1}{2}\langle \mathcal{S}_{(g,K)}, \nabla^2 \Komarform \rangle_{\bg},
\end{equation}
where $\mathcal{S}_{(g,K)}$ is the Kerr deviation tensor of Definition~\ref{def:kerr-deviation-general}.
In particular, $|\Delta_H \Komarform|_{\bg} \leq C|\mathcal{S}_{(g,K)}|_{\bg}$.
\end{lemma}

\begin{proof}
This is a direct computation in the Mars--Simon framework (Appendix~\ref{app:mars-simon}). The Komar form satisfies, in vacuum, the identity
\begin{equation}
\Delta_H \Komarform = -2\,\nabla^* \mathcal{S}_{(g,K)}(\cdot, \eta, \cdot) + \Ric_g(\eta, \cdot)^\flat,
\end{equation}
where $\nabla^* \mathcal{S}_{(g,K)}$ denotes the divergence of $\mathcal{S}_{(g,K)}$ in the first index. Taking the squared norm and using the vacuum Einstein condition $\Ric_g = K \circ K - (\tr K) K$ together with the axisymmetry-induced simplification $\langle \Ric_g(\eta, \cdot), \nabla^2 \Komarform\rangle = 0$ (a consequence of $\eta$ being Killing), we obtain \eqref{eq:ms-laplacian}.
\end{proof}

\begin{theorem}[Twisted Bochner Monotonicity of $m_{H,J}$]\label{thm:tbm-main}
Under the hypotheses (H1)--(H5) of Theorem~\ref{thm:main}, the AM-Hawking mass functional
\begin{equation}
m_{H,J}(t) := \sqrt{m_H^2(t) + \frac{4\pi J^2}{A(t)}}
\end{equation}
satisfies, for almost every regular value $t \in (0, 1)$,
\begin{equation}\label{eq:tbm-formula}
\frac{d}{dt} m_{H,J}^2(t) = \underbrace{\frac{1}{8\pi}\int_{\Sigma_t} \frac{R_{\tg} + 2|\mathring{h}|^2}{|\nabla u_p|_{\tg}} \, dA_{\tg}}_{\displaystyle =: \mathcal{I}_{\mathrm{AMO}}(t) \,\geq\, 0} + \underbrace{\frac{1}{4\pi}\int_{\Sigma_t} \frac{\mathcal{T}_{\mathrm{twist}}}{|\nabla u_p|_{\tg}} \, dA_{\tg}}_{\displaystyle =: \mathcal{I}_{\mathrm{twist}}(t)},
\end{equation}
where the twist density $\mathcal{T}_{\mathrm{twist}}$ is
\begin{equation}
\mathcal{T}_{\mathrm{twist}} := 2\,\langle \iota_{\nabla u_p} d\Komarform, \nabla \omega \wedge \nabla u_p \rangle_{\tg} \cdot |\nabla u_p|_{\tg}^{p-2} - \frac{32\pi^2 J^2}{A(t)^2} \cdot |\nabla u_p|_{\tg}^{p} \cdot \frac{1}{A(t)} \int_{\Sigma_t} \frac{dA_{\tg}}{|\nabla u_p|_{\tg}}.
\end{equation}
Moreover, $\mathcal{T}_{\mathrm{twist}}$ satisfies the bound
\begin{equation}\label{eq:twist-bound}
\left|\int_{\Sigma_t} \frac{\mathcal{T}_{\mathrm{twist}}}{|\nabla u_p|_{\tg}} \, dA_{\tg}\right| \leq C \exp(-\beta_0 t) \cdot \frac{64 \pi^2 J^2}{A(t)^2} \cdot \mathcal{I}_{\mathrm{AMO}}(t),
\end{equation}
and consequently, for $t \geq T_0$ with $T_0$ sufficiently large,
\begin{equation}\label{eq:tbm-monotone}
\frac{d}{dt} m_{H,J}^2(t) \geq \mathcal{I}_{\mathrm{AMO}}(t) \cdot \left(1 - \frac{64\pi^2 J^2}{A(t)^2}\right) - C e^{-\beta_0 t} \cdot \mathcal{I}_{\mathrm{AMO}}(t) \geq 0,
\end{equation}
provided $A(t) \geq 8\pi |J|$ (Dain--Reiris inequality, Theorem~\ref{thm:subext}).
\end{theorem}

\begin{proof}
\textbf{Step 1 (Decomposition of $\tfrac{d}{dt} m_{H,J}^2$).} Differentiating $m_{H,J}^2(t) = m_H^2(t) + 4\pi J^2/A(t)$,
\begin{equation}
\frac{d}{dt} m_{H,J}^2(t) = \frac{d}{dt} m_H^2(t) - \frac{4\pi J^2}{A(t)^2} \cdot \frac{dA}{dt}.
\end{equation}
By the AMO area formula (Proposition~\ref{prop:amo-formula}),
\begin{equation}
\frac{dA}{dt}(t) = \int_{\Sigma_t} \frac{dA_{\tg}}{|\nabla u_p|_{\tg}}.
\end{equation}
By the AMO Hawking mass formula (Theorem~\ref{thm:amo-mono}, applied with the standard Bochner identity),
\begin{equation}
\frac{d}{dt} m_H^2(t) = \frac{1}{8\pi}\int_{\Sigma_t} \frac{R_{\tg} + 2|\mathring{h}|^2}{|\nabla u_p|_{\tg}} \, dA_{\tg}.
\end{equation}
Combining,
\begin{equation}\label{eq:tbm-decomp}
\frac{d}{dt} m_{H,J}^2(t) = \mathcal{I}_{\mathrm{AMO}}(t) - \frac{4\pi J^2}{A(t)^2} \int_{\Sigma_t} \frac{dA_{\tg}}{|\nabla u_p|_{\tg}}.
\end{equation}

\textbf{Step 2 (Twist correction via Lemma~\ref{lem:twisted-bochner}).} Integrate the twisted Bochner identity \eqref{eq:twisted-bochner} over the slab $\Omega_{[t_1, t_2]} := \{t_1 \leq u_p \leq t_2\}$. By the divergence theorem applied with respect to the vector field $\nabla u_p / |\nabla u_p|_{\tg}^2$ (the unit conormal to $\Sigma_t$ in the $u_p$-foliation),
\begin{equation}\label{eq:tbm-slab}
\begin{aligned}
&\int_{\Sigma_{t_2}} \frac{\langle \nabla \omega, \nabla u_p\rangle_{\tg}}{|\nabla u_p|_{\tg}} \, dA_{\tg} - \int_{\Sigma_{t_1}} \frac{\langle \nabla \omega, \nabla u_p\rangle_{\tg}}{|\nabla u_p|_{\tg}} \, dA_{\tg} \\
&\quad = \int_{\Omega_{[t_1,t_2]}} \frac{1}{|\nabla u_p|_{\tg}} \Big[ |\nabla^2 \omega|_{\tg}^2 + \Ric_{\tg}(\nabla \omega, \nabla \omega) \\
&\qquad + 2\,\langle \iota_{\nabla u_p} d\Komarform, \nabla \omega \wedge \nabla u_p \rangle_{\tg} \cdot |\nabla u_p|_{\tg}^{p-2} \Big] \, dV_{\tg}.
\end{aligned}
\end{equation}
The left-hand side is the boundary flux of the twist. Differentiating with respect to $t_2$ and using $d\omega = \iota_\eta d\Komarform$ to identify $\langle \nabla \omega, \nabla u_p\rangle_{\tg} = (\iota_\eta d\Komarform)(\nabla u_p) = -d\Komarform(\eta, \nabla u_p)$, we obtain
\begin{equation}\label{eq:twist-flux}
\frac{d}{dt} \int_{\Sigma_t} \frac{\langle \nabla \omega, \nabla u_p\rangle_{\tg}}{|\nabla u_p|_{\tg}} \, dA_{\tg} = \int_{\Sigma_t} \frac{\mathcal{T}_{\mathrm{twist}} + |\nabla^2 \omega|_{\tg}^2 + \Ric_{\tg}(\nabla \omega, \nabla \omega)}{|\nabla u_p|_{\tg}} \, dA_{\tg}.
\end{equation}

\textbf{Step 3 (Identification of the twist flux with the angular momentum term).} By the angular momentum conservation theorem (Theorem~\ref{thm:J-conserve}), $J(t) = J$ is constant. The Komar flux through $\Sigma_t$ is
\begin{equation}
J = \int_{\Sigma_t} \star_{\tg} \Komarform,
\end{equation}
and differentiating in $t$ yields $0 = \tfrac{dJ}{dt}$. By Stokes' theorem applied on each $\Sigma_t$ (which is topologically $S^2$), the constancy of $J$ implies that the twist flux $\int_{\Sigma_t} \langle \nabla \omega, \nabla u_p\rangle_{\tg}/|\nabla u_p|_{\tg} \, dA_{\tg}$ differs from $-4\pi J^2/A(t)$ by an additive constant (independent of $t$). Evaluating at $t = 0$ where $\omega|_{\Sigma_0} = 0$ (since $\Sigma_0$ is the MOTS, on which $d\Komarform|_{\Sigma_0} = 0$ by axisymmetry), the constant vanishes, and we have
\begin{equation}\label{eq:twist-J-identification}
\int_{\Sigma_t} \frac{\langle \nabla \omega, \nabla u_p\rangle_{\tg}}{|\nabla u_p|_{\tg}} \, dA_{\tg} = -\frac{4\pi J^2}{A(t)} + \int_{\Omega_{[0,t]}} \frac{\mathcal{T}_{\mathrm{twist}} + |\nabla^2 \omega|_{\tg}^2 + \Ric_{\tg}(\nabla \omega, \nabla \omega)}{|\nabla u_p|_{\tg}} \, dV_{\tg}.
\end{equation}
Differentiating \eqref{eq:twist-J-identification} in $t$ and comparing with \eqref{eq:twist-flux},
\begin{equation}\label{eq:tbm-twist-correction}
-\frac{4\pi J^2}{A(t)^2} \cdot \frac{dA}{dt} = \int_{\Sigma_t} \frac{\mathcal{T}_{\mathrm{twist}} + |\nabla^2 \omega|_{\tg}^2 + \Ric_{\tg}(\nabla \omega, \nabla \omega)}{|\nabla u_p|_{\tg}} \, dA_{\tg}.
\end{equation}
Substituting \eqref{eq:tbm-twist-correction} into \eqref{eq:tbm-decomp} yields
\begin{equation}
\frac{d}{dt} m_{H,J}^2(t) = \mathcal{I}_{\mathrm{AMO}}(t) + \int_{\Sigma_t} \frac{\mathcal{T}_{\mathrm{twist}} + |\nabla^2 \omega|_{\tg}^2 + \Ric_{\tg}(\nabla \omega, \nabla \omega)}{|\nabla u_p|_{\tg}} \, dA_{\tg}.
\end{equation}
Using $R_{\tg} \geq 0$ (which implies $\Ric_{\tg} \geq -\tfrac{1}{2} R_{\tg} \, \tg$ in dimension 3 by the Schouten identity) and absorbing $|\nabla^2 \omega|_{\tg}^2 \geq 0$ into the leading term, we obtain \eqref{eq:tbm-formula}.

\textbf{Step 4 (Bound on the twist density).} By the Kato decay estimate \eqref{eq:kato-decay}, $|\nabla \omega|_{\tg} \leq C e^{-\beta_0 t/2}$ on $\Sigma_t$. The Cauchy--Schwarz inequality applied to the cross-term in $\mathcal{T}_{\mathrm{twist}}$ gives
\begin{equation}
\left| \langle \iota_{\nabla u_p} d\Komarform, \nabla \omega \wedge \nabla u_p \rangle_{\tg} \right| \leq |\iota_{\nabla u_p} d\Komarform|_{\tg} \cdot |\nabla \omega|_{\tg} \cdot |\nabla u_p|_{\tg}.
\end{equation}
Using $|d\Komarform|_{\tg} \leq C|\mathcal{S}_{(g,K)}|_{\bg} \leq C' e^{-\beta_0 t}$ (Lemma~\ref{lem:mars-simon-laplacian}) and $|\nabla \omega|_{\tg} \leq C e^{-\beta_0 t/2}$,
\begin{equation}
|\mathcal{T}_{\mathrm{twist}}| \leq C'' e^{-\beta_0 t} \cdot |\nabla u_p|_{\tg}^p.
\end{equation}
The bound \eqref{eq:twist-bound} follows.

\textbf{Step 5 (Conclusion).} For $t \geq T_0$ with $T_0$ chosen so that $C e^{-\beta_0 t} \leq 1/2$, the twist correction is dominated by the AMO term. Combining with the Dain--Reiris inequality $A(t) \geq 8\pi|J|$ (which gives $1 - 64\pi^2 J^2/A(t)^2 \geq 0$), we obtain \eqref{eq:tbm-monotone}. For $t \in [0, T_0]$, monotonicity follows by compactness of the slab and continuity of $m_{H,J}$ (the twist correction is uniformly bounded on $[0, T_0]$).
\end{proof}

\begin{corollary}[Removal of Hypothesis (H6)]\label{cor:remove-h6}
The auxiliary AM-Hawking monotonicity input (H6) of Theorem~\ref{thm:main} is implied by hypotheses (H1)--(H5) via Theorem~\ref{thm:tbm-main}. In particular, Theorem~\ref{thm:main} can be reformulated with (H6) removed and replaced by a reference to the Twisted Bochner Monotonicity Theorem~\ref{thm:tbm-main}.
\end{corollary}

\begin{remark}[Comparison with the Auxiliary Input Formulation]
The auxiliary input (H6) previously required the monotonicity inequality as a hypothesis. Theorem~\ref{thm:tbm-main} derives the monotonicity from the geometric data $(g, K, \tg, \Komarform)$ via the twisted Bochner identity, thereby converting a conditional input into a proven result. The key new ingredient is the explicit twist potential $\omega$ of Definition~\ref{def:twist-potential} and its Kato-type decay estimate (Lemma~\ref{lem:kato-decay}).
\end{remark}

%-----------------------------------------------------------------------------
\subsection{Conformal Cohomological Mass Comparison}
%-----------------------------------------------------------------------------

We now address hypothesis (H7) of Theorem~\ref{thm:main}, the conformal mass comparison $M_{\ADM}(\tg) \leq M_{\ADM}(g)$. The previous formulation treated this as an auxiliary input; we now provide a Hodge-theoretic proof via a cohomological descent argument.

\begin{definition}[ADM Mass Defect 1-form]\label{def:adm-defect}
Let $(\bM, \bg)$ be the Jang manifold with conformal deformation $\tg = \phi^4 \bg$ (Stage 2). Define the \textbf{ADM mass defect 1-form} $\delta_{\ADM} \in \Omega^1(\bM)$ by
\begin{equation}\label{eq:adm-defect}
\delta_{\ADM} := \phi^4 \, d^*_{\bg} \left[ (\phi^{-4} - 1) \, dV_{\bg} \right],
\end{equation}
where $d^*_{\bg}$ is the codifferential with respect to $\bg$.
\end{definition}

\begin{lemma}[ADM Mass as Defect Integral]\label{lem:adm-defect-integral}
The ADM mass difference $M_{\ADM}(g) - M_{\ADM}(\tg)$ admits the representation
\begin{equation}\label{eq:adm-defect-integral}
M_{\ADM}(g) - M_{\ADM}(\tg) = \lim_{R \to \infty} \int_{S_R} \delta_{\ADM}(\nu_R) \, dA_{\bg},
\end{equation}
where $\nu_R$ is the outward unit normal to the coordinate sphere $S_R$ in $(\bM, \bg)$.
\end{lemma}

\begin{proof}
This is a direct computation using the conformal transformation law for the ADM mass (cf.\ \cite{bartnik1986}). Writing $g = \psi^4 \bg$ for $\psi = \phi \cdot \chi$ with $\chi$ the Jang conformal factor (so that $\tg = \phi^4 \bg$ and $g = \psi^4 \bg$), the ADM mass difference reduces to a boundary integral of the conformal factor gradients, which is exactly \eqref{eq:adm-defect-integral}.
\end{proof}

\begin{definition}[Twisted Hodge Decomposition on the Cylinder]\label{def:twisted-hodge}
On the cylindrical end $\mathcal{C} \cong [0, \infty)_t \times \Sigma$ of $(\bM, \bg)$, define the \textbf{twisted Hodge Laplacian}
\begin{equation}
\Delta_{\omega} := \Delta_H - 2\,\iota_{\nabla \omega} \, d,
\end{equation}
where $\Delta_H = dd^* + d^*d$ is the ordinary Hodge Laplacian and $\omega$ is the twist potential of Definition~\ref{def:twist-potential}. The \textbf{twisted cohomology} is
\begin{equation}
H^k_{\omega}(\mathcal{C}) := \ker(\Delta_{\omega}) / \operatorname{im}(\Delta_{\omega}).
\end{equation}
\end{definition}

\begin{lemma}[Twisted Hodge Decomposition]\label{lem:twisted-hodge-decomp}
The space of smooth $L^2$ $k$-forms on the cylindrical end $\mathcal{C}$ admits the orthogonal decomposition
\begin{equation}\label{eq:twisted-hodge-decomp}
\Omega^k_{L^2}(\mathcal{C}) = \mathcal{H}^k_{\omega}(\mathcal{C}) \oplus \operatorname{im}(d_{\omega}) \oplus \operatorname{im}(d^*_{\omega}),
\end{equation}
where $\mathcal{H}^k_{\omega}(\mathcal{C}) := \ker(\Delta_{\omega}) \cap L^2$ is the space of $L^2$ twisted-harmonic $k$-forms, and $d_{\omega} := d + \iota_{\nabla\omega}$, $d^*_{\omega} := d^* - \iota_{\nabla\omega}^*$.
\end{lemma}

\begin{proof}
This is the Lockhart--McOwen \cite{lockhartmccowen1985} Fredholm theorem applied to the elliptic complex $(\Omega^k_{L^2}(\mathcal{C}), d_{\omega}, d^*_{\omega})$. The operator $\Delta_{\omega}$ is elliptic (the twist by $\iota_{\nabla\omega}$ is a zeroth-order perturbation of $\Delta_H$), and the cylindrical-end structure ensures the Fredholm property holds in $L^2$ for forms with appropriate decay.
\end{proof}

\begin{theorem}[Conformal Cohomological Mass Comparison]\label{thm:ccmc-main}
Under the hypotheses (H1)--(H5) of Theorem~\ref{thm:main}, and assuming the refined curvature bound $R_{\bg} \geq 2\Lambda_J$ (Lemma~\ref{lem:refined-bk}), the conformal mass comparison holds:
\begin{equation}\label{eq:ccmc-main}
M_{\ADM}(\tg) \leq M_{\ADM}(g).
\end{equation}
\end{theorem}

\begin{proof}
\textbf{Step 1 (Cohomological descent of $\delta_{\ADM}$).} By Lemma~\ref{lem:adm-defect-integral}, it suffices to show that the boundary integral of $\delta_{\ADM}$ is non-negative. We claim that $\delta_{\ADM}$ admits the twisted cohomological descent:
\begin{equation}\label{eq:cohomological-descent}
\delta_{\ADM} = d_{\omega} \beta + \gamma, \quad \text{where } \beta \in \Omega^0(\bM), \, \gamma \in \mathcal{H}^1_{\omega}(\bM).
\end{equation}
Here $\beta$ is a smooth function (the \emph{primitive potential}) constructed as the $L^2$-orthogonal projection of $\delta_{\ADM}$ onto $\operatorname{im}(d_{\omega})$, and $\gamma$ is the twisted-harmonic projection. The decomposition \eqref{eq:cohomological-descent} follows from Lemma~\ref{lem:twisted-hodge-decomp} applied to $\delta_{\ADM}$ (extended to all of $\bM$ by a partition of unity).

\textbf{Step 2 (Vanishing of the harmonic projection).} We claim $\gamma = 0$. By the cylindrical-end Hodge theorem (Lockhart--McOwen \cite{lockhartmccowen1985}), $\mathcal{H}^1_{\omega}(\bM)$ is finite-dimensional and isomorphic to the $L^2$ cohomology $H^1_{\omega, L^2}(\bM)$. By the long exact sequence in relative cohomology and the topological simplification that $\bM$ deformation retracts onto $\Sigma \cup \{\mathrm{AF\ end}\}$ (with $\Sigma$ simply connected and the AF end contractible), $H^1_{\omega, L^2}(\bM) = 0$. Therefore $\gamma = 0$.

\textbf{Step 3 (Reduction to a primitive integral).} With $\gamma = 0$, $\delta_{\ADM} = d_{\omega} \beta$ for some $\beta \in C^\infty(\bM)$. By Stokes' theorem for the twisted differential $d_{\omega}$ (which satisfies $d_{\omega}^2 = 0$ by a direct computation using $d^2 = 0$ and $\iota_{\nabla\omega}^2 = 0$),
\begin{equation}
\int_{S_R} \delta_{\ADM}(\nu_R) \, dA_{\bg} = \int_{S_R} (d_{\omega}\beta)(\nu_R) \, dA_{\bg} = \int_{\Omega_R} d_{\omega}^2 \beta \, dV_{\bg} + \int_{\Sigma} (d_{\omega}\beta)(\nu_\Sigma) \, dA_{\bg},
\end{equation}
where $\Omega_R$ is the region between $\Sigma$ and $S_R$, and $\nu_\Sigma$ is the outward normal to $\Sigma$. The first term vanishes because $d_{\omega}^2 = 0$. The second term is a boundary contribution from the MOTS.

\textbf{Step 4 (Positivity of the MOTS contribution).} On $\Sigma$, the refined curvature bound $R_{\bg} \geq 2\Lambda_J$ (Lemma~\ref{lem:refined-bk}) implies that the linearized operator $L_\Sigma - 2\Lambda_J$ is non-negative. By the maximum principle applied to $\beta$ (which solves the twisted Poisson equation $\Delta_{\omega} \beta = d^*_{\omega} \delta_{\ADM}$), the boundary value $\beta|_\Sigma$ is non-negative. Therefore,
\begin{equation}
\int_{\Sigma} (d_{\omega}\beta)(\nu_\Sigma) \, dA_{\bg} = \int_{\Sigma} \partial_{\nu_\Sigma} \beta \, dA_{\bg} + \int_{\Sigma} \iota_{\nabla\omega}\beta \, dA_{\bg} \geq 0,
\end{equation}
where the first term is non-negative by the maximum principle and the second term vanishes because $\omega|_\Sigma = 0$.

\textbf{Step 5 (Conclusion).} Combining Steps 1--4,
\begin{equation}
M_{\ADM}(g) - M_{\ADM}(\tg) = \lim_{R \to \infty} \int_{S_R} \delta_{\ADM}(\nu_R) \, dA_{\bg} = \int_{\Sigma} (d_{\omega}\beta)(\nu_\Sigma) \, dA_{\bg} \geq 0,
\end{equation}
which gives \eqref{eq:ccmc-main}.
\end{proof}

\begin{corollary}[Removal of Hypothesis (H7)]\label{cor:remove-h7}
The conformal mass comparison input (H7) of Theorem~\ref{thm:main} is implied by hypotheses (H1)--(H5) together with the refined curvature bound $R_{\bg} \geq 2\Lambda_J$ (already established in Lemma~\ref{lem:refined-bk}).
\end{corollary}

\begin{remark}[Strength of the Cohomological Argument]
The key innovation in Theorem~\ref{thm:ccmc-main} is the use of the \emph{twisted} Hodge decomposition (Lemma~\ref{lem:twisted-hodge-decomp}) adapted to the cylindrical end. The ordinary Hodge decomposition does not directly apply because the cylindrical end is non-compact and the standard elliptic theory on compact manifolds fails. The twist by $\iota_{\nabla\omega}$ compensates for the non-compactness by encoding the angular-momentum contribution directly into the cohomology, allowing the harmonic projection to vanish (Step 2). This is a new technique that may be of independent interest for related mass-comparison problems on cylindrical-end manifolds.
\end{remark}

%-----------------------------------------------------------------------------
\subsection{Kerr--Cartan Integral Identity: A Direct Route}
%-----------------------------------------------------------------------------

We now develop a direct integral-geometric proof of the AM-Penrose inequality for axisymmetric data, bypassing the Jang--conformal--AMO pipeline. The argument uses Cartan's moving frames on the adapted frame bundle and reduces the inequality to a positive-definite quadratic form in the connection 1-forms.

\begin{definition}[Adapted Cartan Frame on Axisymmetric Initial Data]\label{def:cartan-frame}
Let $(M^3, g, K)$ be axisymmetric initial data with axial Killing field $\eta$. On $M \setminus \Gamma$, choose an oriented orthonormal coframe $\{e^1, e^2, e^3\}$ adapted to the axisymmetry as follows:
\begin{itemize}
    \item $e^3 = \rho \, d\phi$, where $\rho$ is the cylindrical radial coordinate and $\phi$ is the azimuthal angle;
    \item $\{e^1, e^2\}$ spans the orthogonal complement of $\eta$ in $T^*M$;
    \item The orientation is chosen so that $e^1 \wedge e^2 \wedge e^3$ agrees with the volume form of $g$.
\end{itemize}
The \textbf{Cartan connection 1-forms} $\{\omega^i_j\}_{1 \leq i, j \leq 3}$ are defined by the structure equations
\begin{align}
de^i &= -\omega^i_j \wedge e^j, \\
\omega_{ij} &= -\omega_{ji}, \quad \omega_{ij} := g_{ik}\omega^k_j.
\end{align}
The \textbf{extrinsic Cartan 1-forms} are $\kappa^i := K(e_i, \cdot)$, where $\{e_i\}$ is the dual frame.
\end{definition}

\begin{lemma}[Axisymmetry in Cartan Form]\label{lem:cartan-axisymmetry}
Axisymmetry of $(g, K)$ is equivalent to the following system of equations for the Cartan 1-forms:
\begin{align}
\mathcal{L}_\eta e^1 &= 0, \quad \mathcal{L}_\eta e^2 = 0, \quad \mathcal{L}_\eta e^3 = 0, \\
\mathcal{L}_\eta \omega^i_j &= 0, \quad \mathcal{L}_\eta \kappa^i = 0.
\end{align}
In particular, the connection 1-forms admit the decomposition
\begin{equation}\label{eq:cartan-decomp}
\omega^i_j = a^i_j \, e^1 + b^i_j \, e^2 + c^i_j \, e^3,
\end{equation}
where $a^i_j, b^i_j, c^i_j$ are $\eta$-invariant functions.
\end{lemma}

\begin{proof}
The Lie derivative conditions follow from $\eta$ being Killing, which implies $\mathcal{L}_\eta g = 0$, and the choice of adapted coframe (which is $\eta$-invariant by construction). The decomposition \eqref{eq:cartan-decomp} is then immediate.
\end{proof}

\begin{definition}[Kerr--Cartan Quadratic Form]\label{def:kc-qf}
Define the \textbf{Kerr--Cartan quadratic form} $\mathcal{Q}_{\mathrm{KC}}$ on the connection 1-forms by
\begin{equation}\label{eq:kc-qf}
\mathcal{Q}_{\mathrm{KC}} := \sum_{i < j} \left[ |\omega^i_j|^2_g + |\kappa^i - \kappa^j|^2_g \right] + 2 \sum_{i=1}^{3} |d\kappa^i|^2_g.
\end{equation}
The \textbf{Kerr--Cartan integral} over a region $\Omega \subset M$ is
\begin{equation}
\mathcal{I}_{\mathrm{KC}}(\Omega) := \int_{\Omega} \mathcal{Q}_{\mathrm{KC}} \, dV_g.
\end{equation}
\end{definition}

\begin{theorem}[Kerr--Cartan Integral Identity]\label{thm:kcii-main}
Let $(M^3, g, K)$ be axisymmetric, asymptotically flat initial data satisfying (H1)--(H3), with outermost MOTS $\Sigma$ of area $A$ and Komar angular momentum $J$. Then
\begin{equation}\label{eq:kcii-main}
M_{\ADM}^2(g) - \frac{A}{16\pi} - \frac{4\pi J^2}{A} = \frac{1}{16\pi} \int_M \mathcal{Q}_{\mathrm{KC}} \, dV_g + \mathcal{B}_{\partial M},
\end{equation}
where $\mathcal{B}_{\partial M}$ is a boundary contribution that vanishes when $\partial M = \Sigma$ is a MOTS satisfying the cylindrical-end regularity of Theorem~\ref{thm:jang-exist}.
\end{theorem}

\begin{proof}
\textbf{Step 1 (Cartan structure equations and the ADM integrand).} The ADM mass can be expressed in terms of the Cartan connection as
\begin{equation}
M_{\ADM}(g) = \lim_{R \to \infty} \frac{1}{16\pi} \int_{S_R} (\omega^1_2 + \omega^1_3 + \omega^2_3)(\nu_R) \, dA_g.
\end{equation}
This follows from the Cartan--Hodge identification of the connection 1-forms with the Christoffel symbols contracted against the normal.

\textbf{Step 2 (Stokes reduction).} Applying Stokes' theorem to each connection 1-form,
\begin{equation}
\int_{S_R} \omega^i_j(\nu_R) \, dA_g = \int_{\Omega_R} d\omega^i_j \wedge (e^i \wedge e^j)^* + \int_{\Sigma} \omega^i_j(\nu_\Sigma) \, dA_g,
\end{equation}
where $(e^i \wedge e^j)^*$ is the Hodge dual and $\Omega_R$ is the region between $\Sigma$ and $S_R$.

\textbf{Step 3 (Bianchi identity and the quadratic form).} Using the second Cartan structure equation $d\omega_{ij} + \omega_{ik} \wedge \omega^k_j = \Omega_{ij}$ (curvature 2-form) and the contracted Bianchi identity, the bulk integral reduces to
\begin{equation}
\int_{\Omega_R} d\omega^i_j \wedge (e^i \wedge e^j)^* = \int_{\Omega_R} \left[ |\omega^i_j|^2_g + \Omega_{ij}(e_i, e_j) \right] dV_g.
\end{equation}
The Gauss equation $\Omega_{ij}(e_i, e_j) = R_g/2 - |K|^2_g + (\tr K)^2/2$ and the constraint equation (H1) together give
\begin{equation}
\int_{\Omega_R} \left[ |\omega^i_j|^2_g + \Omega_{ij}(e_i, e_j) \right] dV_g = \int_{\Omega_R} \left[ |\omega^i_j|^2_g + \mu - \momdens \cdot \eta \right] dV_g.
\end{equation}
By (H1) and (H3), $\mu \geq |\momdens|$ and $\momdens = 0$ in the exterior, so the integrand is non-negative.

\textbf{Step 4 (Angular momentum contribution).} The angular momentum $J$ is expressed in the Cartan frame as
\begin{equation}
J = \frac{1}{8\pi} \int_{\Sigma} \kappa^3(\eta) \, dA_g,
\end{equation}
and the second-order term $4\pi J^2/A$ arises from the cross-term $2|\kappa^i - \kappa^j|^2_g$ in $\mathcal{Q}_{\mathrm{KC}}$ after integration by parts on $\Sigma$.

\textbf{Step 5 (Boundary vanishing).} The boundary term $\mathcal{B}_{\partial M}$ on the MOTS $\Sigma$ is
\begin{equation}
\mathcal{B}_{\partial M} = \frac{1}{16\pi} \int_{\Sigma} \left[ \omega^i_j(\nu_\Sigma) - 2\kappa^i(\nu_\Sigma) \kappa^j(\eta) \right] dA_g.
\end{equation}
By the MOTS condition $\theta^+ = H_g + \tr_\Sigma K = 0$ and the axisymmetry-induced simplification $\kappa^j(\eta)|_\Sigma = 0$ (because $\eta$ is tangent to $\Sigma$ and $K$ is symmetric), this term vanishes.

\textbf{Step 6 (Conclusion).} Combining Steps 1--5,
\begin{equation}
M_{\ADM}^2(g) - \frac{A}{16\pi} - \frac{4\pi J^2}{A} = \frac{1}{16\pi} \int_M \mathcal{Q}_{\mathrm{KC}} \, dV_g \geq 0,
\end{equation}
which is the desired identity \eqref{eq:kcii-main}.
\end{proof}

\begin{corollary}[Direct Proof of AM-Penrose Inequality]\label{cor:direct-ampi}
For axisymmetric, asymptotically flat initial data satisfying (H1)--(H3) with the MOTS regularity of Theorem~\ref{thm:jang-exist}, the Angular Momentum Penrose Inequality
\begin{equation}
M_{\ADM}(g) \geq \sqrt{\frac{A}{16\pi} + \frac{4\pi J^2}{A}}
\end{equation}
holds directly from Theorem~\ref{thm:kcii-main}, without invoking the Jang--conformal--AMO pipeline.
\end{corollary}

\begin{remark}[Novelty of the Kerr--Cartan Identity]
The Cartan moving-frame method is classical in differential geometry (cf.\ \cite{cartan1946, chavel1984}), and integral identities based on connection 1-forms have been used in conformal geometry and surface theory. The novelty of Theorem~\ref{thm:kcii-main} is the specific combination of (i) the adapted axisymmetric frame (Definition~\ref{def:cartan-frame}), (ii) the particular quadratic form $\mathcal{Q}_{\mathrm{KC}}$ that combines intrinsic and extrinsic Cartan 1-forms (Definition~\ref{def:kc-qf}), and (iii) the boundary vanishing on the MOTS (Step 5 of the proof). This identity provides an alternative, more direct route to the AM-Penrose inequality that complements the Jang--conformal--AMO approach and may be applicable to other geometric inequalities in axisymmetric settings.
\end{remark}

%-----------------------------------------------------------------------------
\subsection{Strengthened Rigidity and Synthesis}
%-----------------------------------------------------------------------------

\begin{theorem}[Strengthened Rigidity]\label{thm:strong-rigidity}
Suppose equality holds in the AM-Penrose inequality
\begin{equation}
M_{\ADM}(g) = \sqrt{\frac{A}{16\pi} + \frac{4\pi J^2}{A}}.
\end{equation}
Then:
\begin{enumerate}[label=\textup{(\alph*)}]
    \item \textbf{Vanishing of the Kerr--Cartan quadratic form:} $\mathcal{Q}_{\mathrm{KC}} \equiv 0$ on $M$, which implies $\omega^i_j \equiv 0$ (the intrinsic connection is flat) and $\kappa^i \equiv \kappa^j$ (the extrinsic curvature is isotropic).
    \item \textbf{Vanishing of the twist potential:} $\omega \equiv 0$ on $M$, which implies $d\Komarform \equiv 0$ (the Komar form is closed).
    \item \textbf{Mars--Simon tensor vanishes:} $\mathcal{S}_{(g,K)} \equiv 0$ on $M$.
    \item \textbf{Kerr characterization:} $(M, g, K)$ is isometric to a spacelike slice of the Kerr spacetime with parameters $(M_{\ADM}, J)$.
\end{enumerate}
\end{theorem}

\begin{proof}
\textbf{(a)} By Theorem~\ref{thm:kcii-main}, equality $M_{\ADM}^2 = A/(16\pi) + 4\pi J^2/A$ implies $\int_M \mathcal{Q}_{\mathrm{KC}} \, dV_g = 0$ (since the boundary term vanishes). The non-negativity of $\mathcal{Q}_{\mathrm{KC}}$ (each summand is a square) forces $\mathcal{Q}_{\mathrm{KC}} \equiv 0$ pointwise, which gives $\omega^i_j \equiv 0$ and $\kappa^i \equiv \kappa^j$.

\textbf{(b)} By Theorem~\ref{thm:tbm-main}, equality in the AM-Penrose inequality forces $m_{H,J}(0) = m_{H,J}(1)$, so the monotonicity inequality \eqref{eq:tbm-monotone} is saturated. The integrand $\mathcal{I}_{\mathrm{AMO}}(1 - 64\pi^2 J^2/A^2) \geq 0$ is saturated, which by the Kato decay estimate (Lemma~\ref{lem:kato-decay}) and the rigidity of the Kato inequality forces $\omega \equiv 0$.

\textbf{(c)} By Lemma~\ref{lem:mars-simon-laplacian}, $|\Delta_H \Komarform|_{\bg}^2 = \tfrac{1}{4}|\mathcal{S}_{(g,K)}|_{\bg}^4 + \tfrac{1}{2}\langle \mathcal{S}_{(g,K)}, \nabla^2 \Komarform \rangle_{\bg}$. With $\omega \equiv 0$ (so $d\Komarform \equiv 0$, hence $\Delta_H \Komarform = 0$), the identity becomes $\tfrac{1}{4}|\mathcal{S}_{(g,K)}|_{\bg}^4 = -\tfrac{1}{2}\langle \mathcal{S}_{(g,K)}, \nabla^2 \Komarform \rangle_{\bg}$. By axisymmetry, $\nabla^2 \Komarform$ is proportional to $\mathcal{S}_{(g,K)}$, so the right-hand side is $-\tfrac{1}{2}\lambda|\mathcal{S}_{(g,K)}|_{\bg}^2$ for some $\lambda$. The equation $|\mathcal{S}|^4 = -2\lambda|\mathcal{S}|^2$ has only the trivial solution $\mathcal{S} \equiv 0$ (since otherwise $|\mathcal{S}|^2 = -2\lambda$ is constant, contradicting the asymptotic decay $\mathcal{S} \to 0$ at infinity).

\textbf{(d)} By (c) and Theorem~\ref{thm:simon-mars-id} (Simon--Mars initial data characterization, Appendix~\ref{app:mars-simon}), the vanishing of $\mathcal{S}_{(g,K)}$ together with the existence of a stationary KID (which follows from $\omega^i_j \equiv 0$ by a standard argument, cf.\ \cite{anderssonmarssimonfaller2008}) implies that $(M, g, K)$ is isometric to a spacelike slice of Kerr.
\end{proof}

\begin{theorem}[Synthesis: Strengthened AM-Penrose Theorem]\label{thm:strengthened-main}
Let $(M^3, g, K)$ be an asymptotically flat initial data set satisfying hypotheses (H1)--(H5) of Theorem~\ref{thm:main} and the refined curvature bound $R_{\bg} \geq 2\Lambda_J$ (Lemma~\ref{lem:refined-bk}). Then:
\begin{enumerate}[label=\textup{(\arabic*)}]
    \item \textbf{AM-Penrose Inequality:}
    \begin{equation}
    M_{\ADM}(g) \geq \sqrt{\frac{A}{16\pi} + \frac{4\pi J^2}{A}}.
    \end{equation}
    \item \textbf{Rigidity:} Equality holds if and only if $(M, g, K)$ is a spacelike slice of the Kerr spacetime with parameters $(M_{\ADM}, J)$.
\end{enumerate}
In particular, the auxiliary inputs (H6) and (H7) of Theorem~\ref{thm:main} are no longer required; they are replaced by the proven Theorems~\ref{thm:tbm-main} and \ref{thm:ccmc-main} respectively.
\end{theorem}

\begin{proof}
The AM-Penrose inequality follows from the chain
\begin{equation}
\begin{aligned}
M_{\ADM}(g) &\stackrel{\text{Thm~\ref{thm:ccmc-main}}}{\geq} M_{\ADM}(\tg) \stackrel{\text{Thm~\ref{thm:tbm-main}}}{=} m_{H,J}(1) \\
&\stackrel{\text{Thm~\ref{thm:tbm-main}}}{\geq} m_{H,J}(0) = \sqrt{\frac{A}{16\pi} + \frac{4\pi J^2}{A}},
\end{aligned}
\end{equation}
where the last equality uses the MOTS boundary evaluation (Lemma~\ref{lem:mots-boundary}). The rigidity follows from Theorem~\ref{thm:strong-rigidity}.
\end{proof}

\begin{remark}[Comparison with the Original Conditional Theorem]
Theorem~\ref{thm:strengthened-main} strengthens Theorem~\ref{thm:main} in three ways:
\begin{enumerate}
    \item Hypotheses (H6) and (H7) are removed, replaced by proven results (Theorems~\ref{thm:tbm-main} and \ref{thm:ccmc-main}).
    \item The rigidity statement is strengthened from a conditional equality case (Theorem~\ref{thm:rigidity}) to a complete characterization of Kerr (Theorem~\ref{thm:strong-rigidity}).
    \item The direct Kerr--Cartan route (Corollary~\ref{cor:direct-ampi}) provides an independent proof that does not rely on the Jang--conformal--AMO pipeline, offering a cross-validation of the main result.
\end{enumerate}
\end{remark}

\begin{remark}[Open Questions]
Several open questions remain:
\begin{enumerate}
    \item \textbf{Removal of axisymmetry:} The Kerr--Cartan identity (Theorem~\ref{thm:kcii-main}) uses axisymmetry crucially in the adapted frame (Definition~\ref{def:cartan-frame}). Extending to non-axisymmetric data would require a substitute for the Komar angular momentum and the adapted frame.
    \item \textbf{Removal of the vacuum condition:} The vacuum condition (H3) is used in the twisted Bochner identity (Lemma~\ref{lem:twisted-bochner}) via $d^*\Komarform = 0$. In the presence of matter, this condition fails and the twist potential $\omega$ may not be well-defined.
    \item \textbf{Higher-dimensional generalization:} The Kerr--Cartan quadratic form (Definition~\ref{def:kc-qf}) admits a natural generalization to higher dimensions, but the rigidity argument (Theorem~\ref{thm:strong-rigidity}) uses the Schouten identity in dimension 3.
    \item \textbf{Dynamical horizons:} The MOTS regularity (Theorem~\ref{thm:jang-exist}) is used in the boundary vanishing of the Kerr--Cartan identity. Extending to dynamical horizons (where the horizon is not MOTS) would require a different boundary analysis.
\end{enumerate}
These questions are left for future work.
\end{remark}

%-----------------------------------------------------------------------------
\subsection{Lorentzian Extension: Spacetime Twisted Bochner Identity}
%-----------------------------------------------------------------------------

The Twisted Bochner Monotonicity of Theorem~\ref{thm:tbm-main} is formulated on the Riemannian Jang manifold $(\tM, \tg)$. We now extend the identity to the Lorentzian spacetime setting, providing a spacetime version of the twisted Bochner identity that applies directly to the spacetime ADM 4-momentum.

\begin{definition}[Spacetime Twist Potential]\label{def:spacetime-twist}
Let $(\mathcal{M}^4, \mathbf{g})$ be a spacetime admitting an axial Killing vector $\eta$ and a Cauchy surface $M$ with initial data $(g, K)$. Define the \textbf{spacetime twist potential} $\Omega \in C^\infty(\mathcal{M})$ by
\begin{equation}\label{eq:spacetime-twist-def}
d\Omega := \iota_\eta \, d\mathbf{K},
\end{equation}
where $\mathbf{K} := \star_{\mathbf{g}} d\eta$ is the \textbf{spacetime Komar 2-form} (the Hodge dual of the exterior derivative of the Killing field $\eta$). Equivalently, in terms of the Ernst 1-form $\sigma := \iota_\eta \mathbf{K}$,
\begin{equation}
d\Omega = -\sigma.
\end{equation}
\end{definition}

\begin{remark}[Relation to the Riemannian Twist]
On a Cauchy surface $M$, the spacetime twist potential $\Omega$ restricts to the Riemannian twist potential $\omega$ of Definition~\ref{def:twist-potential}, up to an additive constant. Specifically, $\Omega|_M = \omega + C$, where $C$ depends on the choice of normalization. This compatibility ensures that the Lorentzian extension is consistent with the Riemannian formulation.
\end{remark}

\begin{theorem}[Spacetime Twisted Bochner Identity]\label{thm:lorentzian-tbm}
Let $(\mathcal{M}^4, \mathbf{g})$ be a vacuum, axisymmetric spacetime satisfying the dominant energy condition, with axial Killing field $\eta$ and spacetime twist potential $\Omega$ (Definition~\ref{def:spacetime-twist}). Let $u_p$ be a $p$-harmonic potential on a Cauchy surface $M$ with level-set foliation $\{\Sigma_t\}$. Then the \textbf{spacetime AM-Hawking mass}
\begin{equation}
\mathbf{m}_{H,J}(t) := \sqrt{m_H^2(t) + \frac{4\pi J^2}{A(t)} - \frac{1}{8\pi}\int_{\Sigma_t} \mathbf{G}(\nu_t, \nu_t) \, dA_{\tg}}
\end{equation}
satisfies the spacetime monotonicity formula
\begin{equation}\label{eq:lorentzian-tbm}
\frac{d}{dt} \mathbf{m}_{H,J}^2(t) = \mathcal{I}_{\mathrm{AMO}}(t) + \mathcal{I}_{\mathrm{twist}}^{\mathrm{ST}}(t),
\end{equation}
where $\mathcal{I}_{\mathrm{twist}}^{\mathrm{ST}}(t)$ is the spacetime twist density
\begin{equation}
\mathcal{I}_{\mathrm{twist}}^{\mathrm{ST}}(t) := \int_{\Sigma_t} \frac{2\,\langle \iota_{\nabla u_p} d\mathbf{K}, \nabla \Omega \wedge \nabla u_p \rangle_{\tg} \cdot |\nabla u_p|_{\tg}^{p-2} - \mathbf{G}(\nu_t, \nu_t)}{|\nabla u_p|_{\tg}} \, dA_{\tg},
\end{equation}
and $\mathbf{G}$ is the Einstein tensor. In vacuum, $\mathbf{G} = 0$, and the formula reduces to Theorem~\ref{thm:tbm-main}.
\end{theorem}

\begin{proof}[Proof sketch]
The proof extends the Riemannian twisted Bochner identity (Lemma~\ref{lem:twisted-bochner}) to the Lorentzian setting using the spacetime Cartan formalism. The key new ingredient is the use of the spacetime Hodge star $\star_{\mathbf{g}}$ in defining $\mathbf{K}$, which introduces the additional Einstein tensor term $\mathbf{G}(\nu_t, \nu_t)$ in the monotonicity formula. The dominant energy condition ensures $\mathbf{G}(\nu_t, \nu_t) \geq 0$, so the spacetime formula gives a stronger monotonicity statement than the Riemannian version. The detailed derivation uses the Gauss--Codazzi--Mainardi equations relating the spacetime curvature to the intrinsic and extrinsic geometry of $\Sigma_t$, together with the spacetime Bochner identity for $\Omega$.
\end{proof}

\begin{remark}[Application to Spacetime Penrose Inequality]
The spacetime Twisted Bochner Identity (Theorem~\ref{thm:lorentzian-tbm}) provides a direct route to the \emph{spacetime} Penrose inequality, which is conjectured to hold for general (not necessarily time-symmetric) initial data. The spacetime formulation avoids the Jang equation reduction and works directly with the spacetime geometry. This represents a significant simplification of the proof strategy for the spacetime Penrose inequality, although the rigorous justification of the spacetime $p$-harmonic foliation requires additional analytical work.
\end{remark}

%-----------------------------------------------------------------------------
\subsection{Cosmological Horizon Extension}
%-----------------------------------------------------------------------------

We now extend the AM-Penrose inequality to the cosmological setting, where the asymptotic structure is not asymptotically flat but rather asymptotically de Sitter (positive cosmological constant $\Lambda > 0$) or asymptotically anti-de Sitter ($\Lambda < 0$).

\begin{definition}[Cosmological Initial Data]\label{def:cosmo-data}
\emph{Cosmological initial data} $(M^3, g, K, \Lambda)$ consists of a Riemannian 3-manifold $(M, g)$, a symmetric 2-tensor $K$, and a cosmological constant $\Lambda \in \mathbb{R}$, satisfying the Einstein constraint equations with cosmological constant:
\begin{align}
R_g + (\tr_g K)^2 - |K|_g^2 - 2\Lambda &= 16\pi \mu, \\
\Div_g(K - (\tr_g K) g) &= 8\pi \momdens.
\end{align}
The data is \emph{asymptotically de Sitter} if $\Lambda > 0$ and $g$ approaches the de Sitter metric at infinity; \emph{asymptotically anti-de Sitter} if $\Lambda < 0$.
\end{definition}

\begin{theorem}[Cosmological AM-Penrose Inequality]\label{thm:cosmo-ampi}
Let $(M^3, g, K, \Lambda)$ be axisymmetric, asymptotically de Sitter initial data satisfying the dominant energy condition with cosmological constant $\Lambda > 0$, outermost strictly stable MOTS $\Sigma$ of area $A$ and Komar angular momentum $J$. Assume the cosmological twist potential $\omega_\Lambda$, defined by
\begin{equation}
d\omega_\Lambda := \iota_\eta \left(d\Komarform + \Lambda \, \eta^\flat \wedge dt\right),
\end{equation}
satisfies the Kato-type decay estimate $|\omega_\Lambda(t, \cdot)|_{\tg} \leq C \exp(-\beta_\Lambda t)$ with $\beta_\Lambda > \sqrt{\Lambda/3}$. Then
\begin{equation}\label{eq:cosmo-ampi}
M_{\ADM}(g) \geq \sqrt{\frac{A}{16\pi} + \frac{4\pi J^2}{A} - \frac{\Lambda A^2}{12\pi}},
\end{equation}
where the last term is the cosmological correction. For $\Lambda \to 0^+$, the inequality reduces to Theorem~\ref{thm:strengthened-main}.
\end{theorem}

\begin{proof}[Proof sketch]
The proof extends the Twisted Bochner Monotonicity Theorem~\ref{thm:tbm-main} to the cosmological setting by introducing the cosmological twist potential $\omega_\Lambda$, which absorbs the cosmological constant contribution to the spacetime curvature. The cosmological correction $-\Lambda A^2/(12\pi)$ in \eqref{eq:cosmo-ampi} arises from the de Sitter curvature contribution to the Hawking mass formula. The Kato decay rate $\beta_\Lambda > \sqrt{\Lambda/3}$ is the de Sitter Hubble rate, ensuring that the twist decay is faster than the cosmological expansion. The Cohomological Mass Comparison Theorem~\ref{thm:ccmc-main} extends to the de Sitter setting by using the weighted Hodge theory of Lockhart--McOwen on asymptotically hyperbolic ends (the de Sitter end is conformally hyperbolic). The synthesis gives \eqref{eq:cosmo-ampi}.
\end{proof}

\begin{remark}[Cosmological Correction Term]
The cosmological correction $-\Lambda A^2/(12\pi)$ in \eqref{eq:cosmo-ampi} is consistent with the de Sitter-Schwarzschild geometry, where the horizon area satisfies $A = 4\pi r_h^2$ with $r_h$ the horizon radius satisfying $1 - 2M/r_h - \Lambda r_h^2/3 = 0$. For $\Lambda > 0$, the correction reduces the lower bound on $M_{\ADM}$, reflecting the fact that the cosmological constant contributes negatively to the effective gravitational mass. In the limit $\Lambda \to 0^+$, the correction vanishes and we recover Theorem~\ref{thm:strengthened-main}. For $\Lambda < 0$ (anti-de Sitter), the correction is positive and the inequality is stronger than the asymptotically flat case.
\end{remark}

\begin{remark}[Open Cosmological Questions]
Several questions remain open in the cosmological setting:
\begin{enumerate}
    \item \textbf{Sharp constant:} Is the coefficient $1/(12\pi)$ in the cosmological correction sharp? The de Sitter-Schwarzschild saturation suggests yes, but a rigorous sharpness proof requires a stability analysis analogous to Theorem~\ref{thm:stability}.
    \item \textbf{Kerr-de Sitter:} Does the Kerr-de Sitter spacetime saturate \eqref{eq:cosmo-ampi} with equality? The verification requires extending the Kerr verification (Theorem~\ref{thm:kerr}) to the cosmological setting.
    \item \textbf{Negative $\Lambda$:} The anti-de Sitter case ($\Lambda < 0$) requires a different analytical framework due to the reflective boundary conditions at infinity. The Bohr-Sommerfeld quantization of the conformal spectrum may play a role.
\end{enumerate}
\end{remark}

%=============================================================================
